%! Function Model Selection and Construction — Unit 1, Lesson 1.7 — Homework
\documentclass[10pt]{article}
\usepackage{precalculus-article}
\usepackage{precalculus-boxes}
\newcommand{\TallMath}[1]{$\displaystyle #1\rule[-1.4em]{0pt}{3.2em}$}

\begin{document}

\pageheader{Unit 1, Lesson 1.7}{Homework}

\vspace{0.12in}

\begin{enumerate}[itemsep=10pt]

\item Take differences until they are constant, then name each model type.

\smallskip
\renewcommand{\arraystretch}{1.4}
\begin{center}
\begin{tabular}{|c|c|c|c|c|c|}
\hline
$x$ & 0 & 1 & 2 & 3 & 4 \\
\hline
$y$ & 11 & 8 & 5 & 2 & $-1$ \\
\hline
1st & --- & \blank{1cm} & \blank{1cm} & \blank{1cm} & \blank{1cm} \\
\hline
\end{tabular}
\qquad
\begin{tabular}{|c|c|c|c|c|c|}
\hline
$x$ & 0 & 1 & 2 & 3 & 4 \\
\hline
$y$ & 1 & 3 & 11 & 31 & 69 \\
\hline
1st & --- & \blank{1cm} & \blank{1cm} & \blank{1cm} & \blank{1cm} \\
\hline
2nd & --- & --- & \blank{1cm} & \blank{1cm} & \blank{1cm} \\
\hline
3rd & --- & --- & --- & \blank{1cm} & \blank{1cm} \\
\hline
\end{tabular}
\end{center}

\smallskip
Left table: \blank{2.6cm} \qquad Right table: \blank{2.6cm}

\item Match each context to its model type. Write a letter in each blank:
      \textbf{L} linear, \textbf{Q} quadratic, \textbf{C} cubic,
      \textbf{R} rational, \textbf{P} piecewise.

\begin{enumerate}[label=(\alph*), itemsep=4pt]
  \item A pool is filled at a steady 12 gallons per minute. \blank{1.2cm}
  \item The amount of paint needed to cover a square wall, as a function of the
        wall's side length. \blank{1.2cm}
  \item A city charges \$0.05 per gallon of water for the first 5{,}000 gallons
        and \$0.09 per gallon after that. \blank{1.2cm}
  \item The time to drive 300 miles, as a function of average speed
        (time $=$ 300/speed). \blank{1.2cm}
  \item The amount of concrete needed to fill a cube-shaped footing, as a
        function of the footing's edge length. \blank{1.2cm}
\end{enumerate}

\item \textbf{A rectangle with no wall to lean on.} A gardener has 100 feet of
      fencing for a rectangular plot and must fence \textbf{all four sides}.

\begin{work}
   2w + 2\ell &= 100 \\
         \ell &= 50 - w \\
         A(w) &= w(50 - w) \\
              &= 50w - w^2 \\
            w &= -\frac{50}{2(-1)} \\
              &= 25
\end{work}

\begin{enumerate}[label=(\alph*), itemsep=4pt]
  \item Constraint: \blank{4cm} \qquad so $\ell = $ \blank{2.8cm}
  \item $A(w) = $ \blank{4cm} \qquad Domain: \blank{1cm} $< w <$ \blank{1cm}
  \item The largest area occurs at $w = $ \blank{1.6cm} ft, which makes the plot
        a \blank{2.4cm}. Explain in one sentence why this answer is a square but
        the barn pen from class was not.

        \writeline
\end{enumerate}

\boxguard[24]
\item \textbf{Back to the pen.} The garden club's model is
      $A(w) = 80w - 2w^2$ on $0 < w < 40$.

\begin{work}
   A(5) &= 80(5) - 2(5)^2 \\
        &= 400 - 50 \\
        &= 350 \\
  A(35) &= 80(35) - 2(35)^2 \\
        &= 2800 - 2450 \\
        &= 350 \\
  \frac{A(15) - A(5)}{15 - 5} &= \frac{750 - 350}{10} \\
                              &= 40
\end{work}

\begin{enumerate}[label=(\alph*), itemsep=4pt]
  \item $A(5) = $ \blank{1.8cm} ft$^2$ \qquad $A(35) = $ \blank{1.8cm} ft$^2$.
        Two different widths, the same area --- why does that make sense?

        \blank{7cm}
  \item Average rate of change from $w = 5$ to $w = 15$: \blank{1.8cm},
        with units \blank{5cm}.
  \item Is $w = 45$ in the domain of this model? \blank{1.6cm} \quad Why not?
        \blank{5.6cm}
\end{enumerate}

\item \textbf{An inverse-proportion model.} The force $F$ between two magnets is
      inversely proportional to the square of the distance $d$ between them, so
      $F = \dfrac{k}{d^2}$. At $d = 3$ cm the force measures 20 newtons.

\begin{work}
    20 &= \frac{k}{3^2} \\
     k &= 180 \\
  F(6) &= \frac{180}{6^2} \\
       &= 5
\end{work}

\begin{enumerate}[label=(\alph*), itemsep=4pt]
  \item $k = $ \blank{2cm} \qquad so $F(d) = $ \blank{3cm}
  \item $F(6) = $ \blank{2cm} newtons
  \item Domain restriction: $d$ \blank{2cm}, because \blank{5cm}
\end{enumerate}

\item A table of equally spaced inputs has second differences that are all
      equal to 4. Which statement is correct?

\begin{enumerate}[label=\Alph*., itemsep=3pt]
  \item Linear, because the differences are constant.
  \item Quadratic, because the second differences are constant.
  \item Degree 5, because the table has five rows.
  \item No polynomial model fits, because the second differences are not zero.
\end{enumerate}

\end{enumerate}

\vspace{0.1in}

\boxguard
\begin{extensionbox}
The barn wall turns out to be only \textbf{30 feet long}. If the pen's length
$\ell$ is more than 30 feet, the part beyond the barn has to be fenced too ---
so the fencing equation changes.

\begin{enumerate}[label=(\alph*), itemsep=4pt]
  \item When $\ell > 30$, the fence must cover $2w + \ell + (\ell - 30) = 80$.
        Solve that for $\ell$: \blank{3.4cm}
  \item That case applies when $w$ is \textbf{less than} \blank{1.6cm} feet.
  \item Write the two area rules. \quad $A(w) = w(55 - w)$ for
        \blank{2.6cm}, \quad and $A(w) = w(80 - 2w)$ for \blank{2.6cm}.
  \item Both rules give the same area at the threshold. Check it:
        \blank{2cm} ft$^2$. Why must a piecewise model agree at the threshold?

        \writeline
\end{enumerate}
\end{extensionbox}

\vspace{0.1in}

\boxguard
\begin{notesbox}{Preview of Lesson 2.1}
Every model you built today came out a polynomial, and the useful one came out
quadratic. Unit 2 opens with \textit{Quadratic Functions Revisited}, and
$A(w) = 80w - 2w^2$ is still on the board: where the vertex formula
$-\frac{b}{2a}$ comes from, what vertex form shows that standard form hides, and
why a parabola is symmetric about its vertex --- which is exactly why $A(5)$ and
$A(35)$ came out equal in problem 4.
\end{notesbox}

\end{document}
